\documentclass[12pt]{article}
\usepackage[T1]{fontenc}
\usepackage[utf8]{inputenc}
\usepackage{lmodern}
\usepackage{textcomp}
\usepackage{enumitem}
\usepackage{geometry}
\geometry{margin=1in}
\begin{document}
\begin{center}
Answer Key - Biology - Undergraduate Level Assessment 10 \
\small (Answers are ordered exactly as in the question paper.)
\end{center}

\section*{Multiple Choice}
\begin{enumerate}[label=\arabic*.]
\item \textbf{Answer:} B
\\ \textit{Options:} A) Respiration; B) Photosynthesis; C) Bacterial decomposition; D) Combustion of fossil fuels
\\ \textit{Note:} Photosynthesis is the process by which plants convert carbon dioxide into organic compounds, thereby removing carbon from the atmosphere rather than contributing to it.
\item \textbf{Answer:} B
\\ \textit{Options:} A) Outer membrane, inner membrane, thylakoid membrane; B) Thylakoid membrane, inner membrane, outer membrane; C) Inner membrane, thylakoid membrane, outer membrane; D) Thylakoid membrane, outer membrane, inner membrane; E) Strama, outer membrane, inner membrane
\\ \textit{Note:} The correct sequence of membranes in the chloroplast is thylakoid membrane, inner membrane, and outer membrane, as the thylakoid membrane is located within the stroma of the chloroplast, which is enclosed by the inner membrane, and the inner membrane itself is surrounded by the outer membrane.
\item \textbf{Answer:} D
\\ \textit{Options:} A) All phenotypic variation is the result of genotypic variation.; B) All genetic variation produces phenotypic variation.; C) All nucleotide variability results in neutral variation.; D) All new alleles are the result of nucleotide variability.
\\ \textit{Note:} The correct answer is based on the principle that nucleotide variability, arising from mutations in the DNA sequence, is the source of new alleles, which contribute to genetic variation within populations.
\item \textbf{Answer:} A
\\ \textit{Options:} A) Smaller cells are more resilient to external damage.; B) Smaller cells have a larger number of organelles.; C) Smaller cells require less energy to function.; D) Smaller cells can multiply faster.; E) Smaller cells fit together more tightly.
\\ \textit{Note:} Smaller cells have a higher surface area-to-volume ratio, which enhances their ability to efficiently exchange materials with their environment and makes them less vulnerable to damage from external factors.
\item \textbf{Answer:} D
\\ \textit{Options:} A) Promoter; B) Inducer; C) Repressor; D) Operator; E) Inhibitor
\\ \textit{Note:} The operator is a short sequence of DNA that acts as a binding site for regulatory proteins, which can either enhance or inhibit the transcription of adjacent genes, thus playing a crucial role in gene regulation.
\end{enumerate}

\section*{True / False}
\begin{enumerate}[label=\arabic*.]
\setcounter{enumi}{5}
\item \textbf{Answer:} True
\\ \textit{Options:} A) True; B) False
\\ \textit{Note:} Calcitonin lowers blood calcium levels by inhibiting osteoclast activity in the bones and reducing renal tubular reabsorption of calcium, counteracting the effects of parathormone, which raises blood calcium levels.
\item \textbf{Answer:} False
\\ \textit{Options:} A) True; B) False
\\ \textit{Note:} The probability of rolling a 5 on two successive rolls of a balanced die is actually 1/36, since the probability of rolling a 5 on a single roll is 1/6, and the events are independent, leading to a combined probability of (1/6) * (1/6) = 1/36.
\item \textbf{Answer:} False
\\ \textit{Options:} A) True; B) False
\\ \textit{Note:} Flowers with strong sweet fragrances are primarily adapted to attract bees, as they rely on scent cues for foraging, while hummingbirds are more attracted to bright colors and nectar availability rather than fragrance.
\item \textbf{Answer:} False
\\ \textit{Options:} A) True; B) False
\\ \textit{Note:} The statement is false because meiosis I results in two haploid cells, but these cells are not yet gametes, as they must undergo meiosis II to produce four haploid gametes.
\item \textbf{Answer:} True
\\ \textit{Options:} A) True; B) False
\\ \textit{Note:} Operant conditioning can effectively assess a dog's sensitivity to different sound frequencies by reinforcing specific responses to stimuli, allowing researchers to gauge the dog's ability to perceive and react to varied auditory cues.
\end{enumerate}

\section*{Long Form Response}
\begin{enumerate}[label=\arabic*.]
\setcounter{enumi}{10}
\item \textbf{Answer:} Removing a bacterial cell wall through lysozyme treatment can lead to the formation of protoplasts, which are bacteria without their cell wall. In an isotonic medium, where the osmotic pressure is balanced with that of the bacterial protoplasm, the bacteria can survive temporarily since there is no net movement of water that would cause lysis. Biochemically, bacterial cell walls are primarily composed of peptidoglycan, a polymer of sugars and amino acids that provides structural integrity. In contrast, the cell walls of eukaryotic plants are made of cellulose, a polysaccharide that offers rigidity and protection. This fundamental difference in composition highlights the diverse strategies that bacteria and plants use to maintain cellular structure and function.
\\ \textit{Note:} correct because it accurately describes how lysozyme treatment removes the bacterial cell wall, leading to protoplast formation, and explains that under isotonic conditions, these protoplasts can survive temporarily while also highlighting the biochemical differences between bacterial cell walls, which are made of peptidoglycan, and plant cell walls, which are primarily composed of cellulose.
\item \textbf{Answer:} Inorganic salts, such as sodium and potassium, play a critical role in maintaining fluid balance and osmotic pressure within the bloodstream; however, recent findings suggest that their concentration does not directly influence blood pressure. Instead, blood pressure regulation is more significantly affected by hormonal control and the overall health of the cardiovascular system. Other solutes present in plasma, including proteins, glucose, and waste products, also contribute to blood viscosity and osmotic pressure, which can indirectly influence blood pressure by affecting blood volume and vascular resistance. Therefore, while inorganic salts are essential for various physiological functions, their concentration alone does not determine blood pressure levels.
\\ \textit{Note:} The answer correctly identifies the critical role of inorganic salts, particularly sodium and potassium, in fluid balance and osmotic pressure while emphasizing that blood pressure regulation is primarily influenced by hormonal mechanisms and cardiovascular health, alongside the contributions of other plasma solutes to blood viscosity.
\end{enumerate}

\vfill
\noindent\textit{End of Answer Key}
\end{document}